\section{Proof of \cref{thm:zap}}
\label{sec:proof-zap}
We prove \cref{thm:zap} in this section.
\begin{proof}
\heading{Complexity.} our $\ZAP$ runs in $\aczero$, since the underlying $\NIZK$ runs in $\aczero$.

\heading{Completeness.}
The completeness of $\ZAP$ follows immediately from that of $\NIZK$ and the fact that $\NIZKChe(\crs_0,\NIZKCon(\crs_0))=1$ for all $\crs_0\in\NIZKGen$ (see \cref{def:vc}). 

\heading{Perfect Soundness.}
Due to the verifiable correlated key generation of $\NIZK$, we have $\crs_0\in\NIZKGen$ or $\crs_1\in\NIZKGen$ for a valid proof $\pi=(\crs_0,\crs_1,\pi_0,\pi_1)$. Hence, the perfect soundness of $\ZAP$ follows immediately from that of $\NIZK$.

\heading{$\aczero$-Witness Indistinguishability.}  
We prove the witness indistinguishability of $\ZAP$ by  a sequence of  games as in \cref{fig:games-zap}.

\begin{figure}[htbp]
\begin{center}
\framebox{
\small
  \begin{tabular}{l}
    \begin{minipage}[t]{0.9\textwidth}
    \underline{$\ZAPProve(\x,\w_0)$, \dbox{\fbox{$\ZAPProve(\x,\w_0,\w_1)$}}, \ddbox{\gbox{$\ZAPProve(\x,\w_1)$}}:} \\
     $\gameg_0$, \fbox{$\gameg_1$}, \dbox{$\gameg_2$}, \gbox{$\gameg_3$}, \ddbox{$\gameg_4$}\\
     
   $(\crs_0,\td_0)\getsr\NIZKTGen$, \dbox{$\crs_0\getsr\NIZKGen$}, \ddbox{$(\crs_0,\td_0)\getsr\NIZKTGen$}\\
    $\crs_1=\NIZKCon(\crs_0)$\\
    $\pi_0\getsr\NIZKProve(\crs_0,\x,\w_0)$, \fbox{$\pi_0\getsr\NIZKProve(\crs_0,\x,\w_1)$}\\ 
     $\pi_1\getsr\NIZKProve(\crs_1,\x,\w_0)$, \gbox{$\pi_1\getsr\NIZKProve(\crs_1,\x,\w_1)$}\\
     Return $\pi=(\crs_0,\crs_1,\pi_0,\pi_1)$
    \end{minipage}
          \end{tabular}
      }
\end{center}
\caption{Modifications on $\ZAPProve$ in the intermediate games.}
\label{fig:games-zap}
\end{figure}

%\paragraph{\bf Game $\gameg_0$.}

%\paragraph{\bf Game $\gameg_1$.} 

%\paragraph{\bf Game $\gameg_2$.}

%\paragraph{\bf Game $\gameg_3$.} 

%\paragraph{\bf Game $\gameg_4$.} 

Let $\advA=\{a_\lambda\}_{\lambda\in\N}\in\aczero$ be an adversary against the witness indistinguishability of $\ZAP$. It receives a proof $\pi$ generated by the (modified) prover in each game as defined in \cref{fig:games-zap}. Below by $\varepsilon_i$ we denote the probability that $a_\lambda$ outputs $1$ in Game $\gameg_i$ for $i=0,\cdots,4$.


\heading{Games $\gameg_0$ and $\gameg_1$.}
$\gameg_0$ is the real game where $a_\lambda$ receives 
$\pi=(\crs_0,\crs_1,\pi_0,\allowbreak\pi_1)\getsr\ZAPProve( \x,\w_0)$.
$\gameg_1$ is the same as $\gameg_0$ except that $\pi_0$ is generated as 
$\pi_0\getsr\NIZKProve(\crs_0,\x,\w_1)$ {instead of}
$\pi_0\getsr\NIZKProve(\crs_0,\x,\w_0)$.
\begin{lemma}
\label{lem:zap1}
$\varepsilon_0=\varepsilon_1$.
\end{lemma}
\begin{proof}
\cref{lem:zap1} follows immediately from the composable zero knowledge of $\NIZK$.
\end{proof}

\heading{Game $\gameg_2$.}
This is the same as $\gameg_1$ except that $\crs_0$ is generated  as $\crs_0\getsr\NIZKGen$ instead of $(\crs_0,\td_0)\getsr\NIZKTGen$.
\begin{lemma}
\label{lem:zap2}
There exists an adversary $\advB_1=\{b_\lambda^1\}_{\lambda\in\N}\in\aczero$ such that $b_\lambda^1$ breaks the composable zero-knowledge of $\NIZK$ with probability
$|\varepsilon_2-\varepsilon_1|$.
%$$| \Pr[\gameg^{a_\lambda}_2\Rightarrow 1] - \Pr[\gameg^{a_\lambda}_1\Rightarrow 1]|.$$
\end{lemma}
\begin{proof}
We build the distinguisher $b_\lambda^1$ as follows.

$b_\lambda^1$ runs as in $\gameg_1$ except that now it takes $\crs_0$ as input from the composable zero-knowledge game of $\NIZK$. $\crs_0$ can be generated as $(\crs_0,\td_0)\getsr\NIZKTGen$ or $\crs_0\getsr\NIZKGen$. When $a_\lambda$ outputs $\beta\in\bin$, $b_\lambda^1$ outputs $\beta$ as well.

If $\crs_0$ is generated as $(\crs_0,\td_0)\getsr\NIZKTGen$ (respectively, $\crs_0\getsr\NIZKGen$), the view of $a_\lambda$ is the same as its view in $\gameg_1$ (respectively, $\gameg_2$). Hence, the probability that $b_\lambda^1$ breaks the fine-grained matrix linear assumption is
$|\varepsilon_2-\varepsilon_1|$.

Moreover, since all operations in $b_\lambda^1$ are performed in $\aczero$, we have $\B_1 = \{b_{\lambda}^1\}_{\lambda\in\N} \in \aczero$, completing this part of the proof.
\end{proof}

\heading{Game $\gameg_3$.}
$\gameg_3$ is the same as $\gameg_2$ except that $\pi_1$ is generated as $\pi_1\getsr\NIZKProve(\crs_1,\allowbreak\x,\w_1)$ instead of $\pi_1\getsr\NIZKProve(\crs_1,\x,\w_0)$.
\begin{lemma}
\label{lem:zap3}
$\varepsilon_3=\varepsilon_2$.
\end{lemma}
\begin{proof}
By the  verifiable correlated key generation, the distribution of $\NIZKCon(\crs_0)$ is the same as $\crs_1$ for $\crs_0\getsr\NIZKGen$ and $(\crs_1,\td_1)\getsr\NIZKTGen$.
Then \cref{lem:zap3} follows from the composable zero-knowledge of $\NIZK$.
\end{proof}

\heading{Game $\gameg_4$.}
$\gameg_4$ is the same as $\gameg_3$ except that $\crs_0$ is generated as 
$(\crs_0,\td_0)\getsr\NIZKTGen$ instead of $\crs_0\getsr\NIZKGen$.

\begin{lemma}
\label{lem:zap4}
There exists an adversary $\advB_2=\{b_\lambda^2\}_{\lambda\in\N}\in\aczero$ such that $b_\lambda^2$ breaks the composable zero-knowledge of $\NIZK$ with probability
$|\varepsilon_4-\varepsilon_3|$.
\end{lemma}
\begin{proof}
We build the distinguisher $b_\lambda^2$ as follows. 

$b_\lambda^2$ runs as in $\gameg_3$ except that $\crs_0$ is taken as input from its composable zero-knowledge challenger, namely, $\crs_0$ can be generated as $\crs_0\getsr\NIZKGen$ or $(\crs_0,\td_0)\getsr\NIZKTGen$.
When $a_\lambda$ outputs $\beta\in\bin$, $b_\lambda^2$ outputs $\beta$ as well.

If $\crs_0$ is generated as $\crs_0\getsr\NIZKGen$ (respectively, $(\crs_0,\td_0)\getsr\NIZKTGen$), the view of $a_\lambda$ is the same as its view in $\gameg_3$ (respectively, $\gameg_4$). Hence, the probability that $b_\lambda^2$ breaks the composable zero-knowledge of $\NIZK$ is
$|\varepsilon_4-\varepsilon_3|$.

Moreover, since all operations in $b_\lambda^2$ are performed in $\aczero$, we have $\B_2 = \{b_{\lambda}^2\}_{\lambda\in\N} \in \aczero$, and this completes the proof.
\end{proof}


Putting all the above together, \cref{thm:zap} immediately follows.
\end{proof}

\heading{Remark on Non-Interactive Zap for $\NP$.}
Similar to the work of Wang and Pan \cite{EC:WanPan22},
our transformation from NIZK to the non-interactive zap also works for polynomial-time provers, namely, we have an unconditionally secure non-interactive zap for all $\NP$ against $\aczero$ adversaries if we allow polynomial-time provers. In our transformation, generating a zap proof (see \cref{fig:con-zap}) involves two proofs of the underlying NIZK. In this case, we have to show that the above reductions run in $\aczero$, i.e., we need to ensure that they can generate proofs of the underlying NIZK in $\aczero$. This is possible for our NIZK in \cref{fig:con-acnizk}. More precisely, to generate a NIZK proof for an $\NP$ statement, $\aczero$-reductions can perform all the steps except for extending the witness (since the commitments and OR-proofs can be generated in parallel). Extending the witness is not necessary, since the extended witness can be hard-wired in an $\aczero$-reduction beforehand, due to the fact that any statement $\x$ and its two witnesses $\w_0$ and $\w_1$ are a-prior fixed in the hybrid games.

%\section{Proof of \cref{thm:ursnizk}}
%\label{sec:proof-urs}
%In this section, we give the proof of \cref{thm:ursnizk}.
%\begin{proof}
%First, we note that $\URSNIZK\in\aczero$, since $\NIZK\in\aczero$ and running $\NIZK$ in parallel does not increase the circuit depth.
%
%\heading{Completeness.}
%The completeness of $\URSNIZK$ follows immediately from the completeness of $\NIZK$.
%
%\heading{Statistical soundness.}
%For each $(\vec r_{ij},\vec s_{ij})_{j=1}^n$ in $\urs$ generated by $\UGen$, the probability that $(\F(\vec r_{ij})||\vec s_{ij})_{j=1}^n\in\NIZKGen$ is $1/2^n$ according to \cref{lem:v}, where $\{\NIZKGen\}_{\lambda\in\N}$ is the family of (binding) CRS generation algorithms of $\NIZK$. Then, for a randomly chosen URS, the probability that $(\F(\vec r_{ij})||\vec s_{ij})_{j=1}^n\notin\NIZKGen$ for all $\ell$ is $(1-1/2^n)^\ell$. Hence, the probability that there exists a valid proof for a statement $\x\notin\lang_\rho$ is at most $(1-1/2^n)^\ell$, due to the perfect soundness of $\NIZK$. Since $n$ is a constant and $\ell$ is a polynomial, $(1-1/2^n)^\ell$ is negligible, completing this part of proof.
%
%\heading{$\aczero$-composable zero-knowledge.} 
%To prove composable zero-knowledge, we define a sequence of hybrid games $\gameg_0$ to $\gameg_{n\cdot l}$. 
%In $\gameg_0$ (respectively, $\gameg_{n\cdot l}$), an $\aczero$ adversary recieves $\urs=((\vec r_{ij},\vec s_{ij})_{j=1}^n)_{i=1}^\ell$ generated by $\UGen$ (respectively, $\UTGen$), and for each $i\in[\ell]$ and $j\in[n]$, $\gameg_{n\cdot (i-1)+j-1}$ is the same as $\gameg_{n\cdot (i-1)+j}$ except that in $\gameg_{n\cdot (i-1)+j-1}$, the distribution of $(\vec r_{ij},\vec s_{ij})$ is uniform in $\bin^{\lambda(\lambda-1)/2}\times\bin^\lambda$. Without loss of generality, we assume that the $j$th matrix in a CRS of $\NIZK$ is sampled by $\ZeroSamp(\lambda)$.
%\begin{lemma}
%\label{lem:urs}
%For each $i\in[\ell]$ and $j\in[n]$, there exists an adversary $\advB_{n\cdot (i-1)+j-1}=\{b_\lambda^{n\cdot (i-1)+j-1}\}_{\lambda\in\N}\in\aczero$ distinguishing $\matr M\getsr\ZeroSamp(\lambda)$ and $\F(\vec r_{ij})||\vec s_{ij}$, where $\vec r_{ij}\getsr\bin^{\lambda(\lambda-1)/2}$ and $\vec s_{ij}\getsr\bin^{\lambda-1}$, with advantage
%$|\Pr[\gameg_{n\cdot (i-1)+j}^{a_\lambda} \Rightarrow 1] -\Pr[\gameg_{n\cdot (i-1)+j-1}^{a_\lambda} \Rightarrow 1]|$.
%%$$| \Pr[\gameg^{a_\lambda}_2\Rightarrow 1] - \Pr[\gameg^{a_\lambda}_1\Rightarrow 1]|.$$
%\end{lemma}
%\begin{proof}
%We build the distinguisher $b_\lambda^{n\cdot (i-1)+j-1}$ as follows. 
%
%$b_\lambda^{n\cdot (i-1)+j-1}$ runs in exactly the same way as the challenger of $\gameg_{{n\cdot (i-1)+j-1}}$ except that instead of generating $\F(\vec r_{ij})||\vec s_{ij}$ by itself, it takes as input $\F(\vec r_{ij})||\vec s_{ij}$ from its own challenger. When $a_\lambda$ outputs $\beta\in\bin$, $b_\lambda^{n\cdot (i-1)+j-1}$ outputs $\beta$ as well.
%
%If $\F(\vec r_{ij})||\vec s_{ij}$ is generated as $\vec r_{ij}\getsr\bin^{\lambda(\lambda-1)/2}$ and $\vec s_{ij}\getsr\bin^\lambda$ (respectively, $\F(\vec r_{ij})||\vec s_{ij}\getsr\ZeroSamp(\lambda)$), the view of $a_\lambda$ is the same as its view in $\gameg_{n\cdot (i-1)+j-1}$ (respectively, $\gameg_{n\cdot (i-1)+j}$). Hence, the advantage of $b_\lambda^{n\cdot (i-1)+j-1}$ is
%$$| \Pr[\gameg^{a_\lambda}_{n\cdot (i-1)+j}\Rightarrow 1] - \Pr[\gameg^{a_\lambda}_{n\cdot (i-1)+j-1} \Rightarrow 1]|.$$
%
%Moreover, since all operations in $b_\lambda^{n\cdot (i-1)+j-1}$ are performed in $\aczero$, we have $\B_{n\cdot (i-1)+j-1} = \{b_{\lambda}^{n\cdot (i-1)+j-1}\}_{\lambda\in\N} \in \aczero$, completing this part of proof.
%\end{proof}
%
%Then, according to \cref{lem:u}, we conclude that $|\Pr[\gameg_{n\cdot (i-1)+j}^{a_\lambda} \Rightarrow 1] -\Pr[\gameg_{n\cdot (i-1)+j-1}^{a_\lambda} \Rightarrow 1]|\leq\negl(\lambda)$.
%
%One the other hand, if the $j$th matrix in a CRS of $\NIZK$ is sampled by $\OneSamp(\lambda)$, then the indistinguishability  between $\gameg_{n\cdot (i-1)+j-1}$ and $\gameg_{n\cdot (i-1)+j}$ can be shown in the same way according to \cref{cor:u}.
%
%Moreover, due to the $\aczero$-composable zero-knowledge of $\NIZK$, the distributions of $\NIZKProve(\crs_i,\x,\w)$ and $\NIZKSim(\crs_i,\td_i,\x)$ are identical, where $\rho\in\dist_\lambda$, $(\x,\w)$ is a valid statement/witness pair, and $(\crs_i,\td_i)\getsr\NIZKTGen$, for all $i\in[\ell]$. Hence, the $\aczero$-composable zero-knowledge of $\URSNIZK$ holds. 
%%Then $\aczero$-composable zero-knowledge follows from the indistinguishability between the URSs in $\gameg_{n\cdot (i-1)+j-1}$ and $\gameg_{n\cdot (i-1)+j}$ for all $i\in[\ell]$ and $j\in[n]$ due to \cref{lem:u}. Notice that the reductions can be performed in $\aczerotwo$ by running $\NIZK$ in parallel.
%
%Putting all the above together, \cref{thm:ursnizk} immediately follows.
%\end{proof}
%




